\documentclass[12pt,aspectratio=1610]{beamer}
\geometry{paperwidth=20cm,paperheight=12.5cm}
\usepackage[italian,provide=*]{babel}
\usepackage{lmodern,fontspec}
\setsansfont{NimbusSans-Regular.otf}[Path=sorgenti/font/,BoldFont=NimbusSans-Bold.otf,ItalicFont=NimbusSans-Italic.otf]
\setmainfont{NimbusSans-Regular.otf}[Path=sorgenti/font/,BoldFont=NimbusSans-Bold.otf,ItalicFont=NimbusSans-Italic.otf]
\setmonofont{lmmono10-regular.otf}[Path=sorgenti/font/,Scale=.88]
\usefonttheme{professionalfonts}
\usepackage{amsmath,amssymb,booktabs,array,tabularx,graphicx,tikz,listings,bookmark}
\usetikzlibrary{calc,arrows.meta,positioning}
\definecolor{Paper}{HTML}{F8F6F1}
\definecolor{Ink}{HTML}{253344}\definecolor{Navy}{HTML}{15283E}
\definecolor{Quantum}{HTML}{5146A5}\definecolor{Classical}{HTML}{087C74}
\definecolor{Measure}{HTML}{A05B10}\definecolor{Rose}{HTML}{9B4564}
\definecolor{Muted}{HTML}{657180}\definecolor{Rule}{HTML}{D6DAD9}
\definecolor{Wash}{HTML}{EAE7F6}\definecolor{Night}{HTML}{112236}
\definecolor{Lilac}{HTML}{ABB5FF}\definecolor{Mint}{HTML}{72D8C8}\definecolor{Gold}{HTML}{EEC184}
\setbeamercolor{background canvas}{bg=Paper}\setbeamercolor{normal text}{fg=Ink}
\setbeamertemplate{navigation symbols}{}\setbeamertemplate{headline}{}\setbeamertemplate{footline}{}
\setbeamersize{text margin left=0pt,text margin right=0pt}\setbeameroption{hide notes}
\newsavebox{\BodyBox}\newsavebox{\SideBox}\newsavebox{\RecallBox}\newsavebox{\TitleBox}
\newcommand{\MeasureBox}[2]{\typeout{GEOBOX:#1:\the\wd#2:\the\ht#2:\the\dp#2}}
\lstset{language=R,basicstyle=\ttfamily\fontsize{10.2}{13.4}\selectfont,
columns=fullflexible,keepspaces=true,showstringspaces=false,breaklines=true,breakatwhitespace=false,
commentstyle=\color{Muted},keywordstyle=\color{Classical},stringstyle=\color{Quantum},
aboveskip=0pt,belowskip=0pt,xleftmargin=0pt,xrightmargin=0pt,tabsize=2}
\hypersetup{colorlinks=true,linkcolor=Classical,urlcolor=Classical,pdfpagemode=UseOutlines,pdfstartview=Fit,bookmarksopen=true,bookmarksopenlevel=0}
\bookmarksetup{open,openlevel=0,depth=2}
\tikzset{arr/.style={-{Stealth[length=2mm]},line width=.8pt,draw=Classical}}
\hypersetup{pdftitle={Campionamento, stima e incertezza},pdfauthor={},pdfsubject={Corso di statistica e data science per Geologia, primo anno; modulo 05}}
\usetikzlibrary{decorations.pathreplacing,fit,backgrounds,matrix}\hypersetup{pdftitle={P06 · Dai campioni alla validazione spaziale},pdfsubject={Revisione integrale R6}}\begin{document}
\begin{frame}[plain]\hypertarget{slide1}{}\begin{tikzpicture}[remember picture,overlay]\fill[Night] (current page.south west) rectangle (current page.north east);\draw[Mint,line width=2pt] ([xshift=1.05cm,yshift=-1.35cm]current page.north west)--([xshift=3cm,yshift=-1.35cm]current page.north west);\node[anchor=north west,inner sep=0] at ([xshift=1.05cm,yshift=-1.7cm]current page.north west){\color{Mint}\bfseries\fontsize{11}{14}\selectfont P06 · REVISIONE R6};
\node[anchor=north west,inner sep=0] at ([xshift=1.05cm,yshift=-2.65cm]current page.north west){\begin{minipage}{15.8cm}\raggedright\color{white}\bfseries\fontsize{31}{38}\selectfont Dai campioni alla validazione spaziale\end{minipage}};
\node[anchor=north west,inner sep=0] at ([xshift=1.05cm,yshift=-6.4cm]current page.north west){\begin{minipage}{15.8cm}\raggedright\color{Lilac}\fontsize{15}{21}\selectfont Vicinato, permutazioni e trasferimento verso un settore escluso\end{minipage}};
\node[anchor=north west,inner sep=0] at ([xshift=1.05cm,yshift=-8.6cm]current page.north west){\begin{minipage}{16.5cm}\color{white}\fontsize{11}{15}\selectfont 17 slide · teoria, derivazioni, figure e applicazioni in R.\par\smallskip Provenienza dei dati e sorgenti nei materiali allegati.\end{minipage}};
\node[anchor=south east,inner sep=0] at ([xshift=18.95cm,yshift=-11.35cm]current page.north west){\color{Lilac!45}\bfseries\fontsize{64}{68}\selectfont P06};
\end{tikzpicture}\end{frame}
\begin{frame}[plain]\hypertarget{slide2}{}\bookmark[level=0,page=2]{Indice}\begin{tikzpicture}[remember picture,overlay]\node[anchor=north west,inner sep=0] at ([xshift=0.95cm,yshift=-0.65cm]current page.north west){\color{Navy}\bfseries\fontsize{22}{27}\selectfont Percorso del modulo};
\draw[Classical,line width=1.3pt] ([xshift=.95cm,yshift=-1.68cm]current page.north west)--([xshift=2.25cm,yshift=-1.68cm]current page.north west);\node[anchor=north west,inner sep=0] at ([xshift=0.95cm,yshift=-2.23cm]current page.north west){\color{Classical}\bfseries\fontsize{13}{16}\selectfont 3};
\node[anchor=north west,inner sep=0] at ([xshift=2.0cm,yshift=-2.23cm]current page.north west){\begin{minipage}{7.25cm}\raggedright\color{Navy}\bfseries\fontsize{12.4}{15.5}\selectfont\hyperlink{slide3}{Geometria e osservazioni sintetiche}\end{minipage}};
\node[anchor=north west,inner sep=0] at ([xshift=0.95cm,yshift=-3.5300000000000002cm]current page.north west){\color{Classical}\bfseries\fontsize{13}{16}\selectfont 7};
\node[anchor=north west,inner sep=0] at ([xshift=2.0cm,yshift=-3.5300000000000002cm]current page.north west){\begin{minipage}{7.25cm}\raggedright\color{Navy}\bfseries\fontsize{12.4}{15.5}\selectfont\hyperlink{slide7}{Vicinato e riferimento di permutazione}\end{minipage}};
\node[anchor=north west,inner sep=0] at ([xshift=10.299999999999999cm,yshift=-2.23cm]current page.north west){\color{Classical}\bfseries\fontsize{13}{16}\selectfont 11};
\node[anchor=north west,inner sep=0] at ([xshift=11.35cm,yshift=-2.23cm]current page.north west){\begin{minipage}{7.25cm}\raggedright\color{Navy}\bfseries\fontsize{12.4}{15.5}\selectfont\hyperlink{slide11}{Previsione e verifica nel settore orientale}\end{minipage}};
\node[anchor=north west,inner sep=0] at ([xshift=0.95cm,yshift=-10.55cm]current page.north west){\begin{minipage}{18cm}\color{Muted}\fontsize{10}{13}\selectfont I numeri indicano le pagine iniziali; i titoli sono cliccabili. Codice, figure e verifiche sono inclusi nel progetto.\end{minipage}};
\end{tikzpicture}\end{frame}
\begin{lrbox}{\TitleBox}\begin{minipage}[t]{18.2cm}\raggedright\sloppy\color{Ink}\fontsize{20.7}{24}\selectfont \bfseries\color{Navy}La domanda collega struttura e trasferimento\end{minipage}\end{lrbox}
\begin{lrbox}{\BodyBox}\begin{minipage}[t]{10.25cm}\raggedright\sloppy\color{Ink}\fontsize{12.4}{16.4}\selectfont Il progetto utilizza sessantaquattro concentrazioni sintetiche collocate su un reticolo locale, per esaminare due problemi distinti: la presenza di associazione fra valori vicini e la capacità di prevedere un settore geografico escluso dall’adattamento. Il primo problema viene formulato attraverso un indice di Moran e un riferimento di permutazione; il secondo richiede previsioni confrontate con valori tenuti separati. Una struttura spaziale evidente può coesistere con errori elevati nel trasferimento, soprattutto quando il settore da prevedere supera il dominio osservato durante l’adattamento.\end{minipage}\end{lrbox}
\begin{lrbox}{\SideBox}\begin{minipage}[t]{7.45cm}\raggedright\sloppy\color{Ink}\fontsize{12.2}{16.1}\selectfont \input{tikz/P06/D01.tex}\end{minipage}\end{lrbox}
\begin{lrbox}{\RecallBox}\begin{minipage}[t]{16.45cm}\raggedright\sloppy\color{Ink}\fontsize{10.1}{13.2}\selectfont Il test di associazione spaziale e la verifica predittiva sostengono conclusioni differenti.\end{minipage}\end{lrbox}
\MeasureBox{3-TitleBox}{\TitleBox}
\MeasureBox{3-BodyBox}{\BodyBox}
\MeasureBox{3-SideBox}{\SideBox}
\MeasureBox{3-RecallBox}{\RecallBox}
\begin{frame}[plain]\hypertarget{slide3}{}
\bookmark[level=0,page=3]{Geometria e osservazioni sintetiche}
\bookmark[level=1,page=3]{La domanda collega struttura e trasferimento}
\begin{tikzpicture}[remember picture,overlay]\node[anchor=north east,inner sep=0] at ([xshift=19.05cm,yshift=-0.16cm]current page.north west){\color{Classical}\fontsize{8}{10}\selectfont\hyperlink{slide2}{Indice}};
\node[anchor=north west,inner sep=0] at ([xshift=0.95cm,yshift=-0.62cm]current page.north west){\usebox{\TitleBox}};
\draw[Classical,line width=1.3pt] ([xshift=.95cm,yshift=-1.68cm]current page.north west)--([xshift=2.25cm,yshift=-1.68cm]current page.north west);\node[anchor=north west,inner sep=0] at ([xshift=0.95cm,yshift=-2.12cm]current page.north west){\usebox{\BodyBox}};
\node[anchor=north west,inner sep=0] at ([xshift=11.6cm,yshift=-2.1cm]current page.north west){\color{Quantum}\bfseries\fontsize{11.5}{14}\selectfont Lettura del diagramma};
\node[anchor=north west,inner sep=0] at ([xshift=11.6cm,yshift=-2.75cm]current page.north west){\usebox{\SideBox}};
\draw[Rule,line width=.5pt] ([xshift=.95cm,yshift=-10.43cm]current page.north west)--([xshift=19.05cm,yshift=-10.43cm]current page.north west);\node[anchor=north west,inner sep=0] at ([xshift=0.95cm,yshift=-10.66cm]current page.north west){\color{Classical}\bfseries\fontsize{9.4}{12}\selectfont Richiami};
\node[anchor=north west,inner sep=0] at ([xshift=2.6cm,yshift=-10.64cm]current page.north west){\usebox{\RecallBox}};
\draw[Rule,line width=.35pt] ([xshift=.95cm,yshift=-11.92cm]current page.north west)--([xshift=19.05cm,yshift=-11.92cm]current page.north west);\node[anchor=north west,inner sep=0] at ([xshift=0.95cm,yshift=-12.06cm]current page.north west){\color{Muted}\fontsize{7.7}{9.3}\selectfont P06 · Dai campioni alla validazione spaziale};
\node[anchor=north east,inner sep=0] at ([xshift=17.65cm,yshift=-12.06cm]current page.north west){\color{Classical}\fontsize{7.3}{9.1}\selectfont };
\node[anchor=north east,inner sep=0] at ([xshift=19.05cm,yshift=-12.06cm]current page.north west){\color{Muted}\fontsize{7.7}{9.3}\selectfont 3/17};
\end{tikzpicture}\end{frame}
\begin{lrbox}{\TitleBox}\begin{minipage}[t]{18.2cm}\raggedright\sloppy\color{Ink}\fontsize{20.7}{24}\selectfont \bfseries\color{Navy}Le coordinate locali delimitano il significato metrico\end{minipage}\end{lrbox}
\begin{lrbox}{\BodyBox}\begin{minipage}[t]{10.25cm}\raggedright\sloppy\color{Ink}\fontsize{12.4}{16.4}\selectfont Le coordinate x e y variano da zero a 1750 metri con passo di 250 metri, formando otto righe e otto colonne. Le coordinate geografiche incluse nella tabella derivano invece da una costruzione approssimata, della quale va conservata la natura didattica. Le distanze del progetto vengono calcolate sul piano locale metrico, senza attribuirgli automaticamente un sistema UTM reale. Un intervallo numerico plausibile di longitudine e latitudine non certifica datum, accuratezza o provenienza delle posizioni.\par\vspace{.28cm}{\centering\fontsize{13}{17}\selectfont\(\displaystyle\begin{gathered}d_{ij}=\sqrt{(x_i-x_j)^2+(y_i-y_j)^2}\end{gathered}\)\par}\end{minipage}\end{lrbox}
\begin{lrbox}{\SideBox}\begin{minipage}[t]{7.45cm}\raggedright\sloppy\color{Ink}\fontsize{12.2}{16.1}\selectfont $(x_i,y_i)$ e $(x_j,y_j)$: posizioni dei siti i e j; $d_{ij}$: distanza in metri. Si controllano identificativi, coordinate finite e assenza di duplicati. Metadati di riferimento e protocollo di acquisizione andrebbero documentati, per osservazioni reali, prima di trasformazioni cartografiche.\end{minipage}\end{lrbox}
\begin{lrbox}{\RecallBox}\begin{minipage}[t]{16.45cm}\raggedright\sloppy\color{Ink}\fontsize{10.1}{13.2}\selectfont La distanza eredita le unità e le ipotesi geometriche del sistema di coordinate utilizzato.\end{minipage}\end{lrbox}
\MeasureBox{4-TitleBox}{\TitleBox}
\MeasureBox{4-BodyBox}{\BodyBox}
\MeasureBox{4-SideBox}{\SideBox}
\MeasureBox{4-RecallBox}{\RecallBox}
\begin{frame}[plain]\hypertarget{slide4}{}
\bookmark[level=1,page=4]{Le coordinate locali delimitano il significato metrico}
\begin{tikzpicture}[remember picture,overlay]\node[anchor=north east,inner sep=0] at ([xshift=19.05cm,yshift=-0.16cm]current page.north west){\color{Classical}\fontsize{8}{10}\selectfont\hyperlink{slide2}{Indice}};
\node[anchor=north west,inner sep=0] at ([xshift=0.95cm,yshift=-0.62cm]current page.north west){\usebox{\TitleBox}};
\draw[Classical,line width=1.3pt] ([xshift=.95cm,yshift=-1.68cm]current page.north west)--([xshift=2.25cm,yshift=-1.68cm]current page.north west);\node[anchor=north west,inner sep=0] at ([xshift=0.95cm,yshift=-2.12cm]current page.north west){\usebox{\BodyBox}};
\node[anchor=north west,inner sep=0] at ([xshift=11.6cm,yshift=-2.1cm]current page.north west){\color{Quantum}\bfseries\fontsize{11.5}{14}\selectfont Sviluppo e interpretazione};
\node[anchor=north west,inner sep=0] at ([xshift=11.6cm,yshift=-2.75cm]current page.north west){\usebox{\SideBox}};
\draw[Rule,line width=.5pt] ([xshift=.95cm,yshift=-10.43cm]current page.north west)--([xshift=19.05cm,yshift=-10.43cm]current page.north west);\node[anchor=north west,inner sep=0] at ([xshift=0.95cm,yshift=-10.66cm]current page.north west){\color{Classical}\bfseries\fontsize{9.4}{12}\selectfont Richiami};
\node[anchor=north west,inner sep=0] at ([xshift=2.6cm,yshift=-10.64cm]current page.north west){\usebox{\RecallBox}};
\draw[Rule,line width=.35pt] ([xshift=.95cm,yshift=-11.92cm]current page.north west)--([xshift=19.05cm,yshift=-11.92cm]current page.north west);\node[anchor=north west,inner sep=0] at ([xshift=0.95cm,yshift=-12.06cm]current page.north west){\color{Muted}\fontsize{7.7}{9.3}\selectfont P06 · Dai campioni alla validazione spaziale};
\node[anchor=north east,inner sep=0] at ([xshift=17.65cm,yshift=-12.06cm]current page.north west){\color{Classical}\fontsize{7.3}{9.1}\selectfont };
\node[anchor=north east,inner sep=0] at ([xshift=19.05cm,yshift=-12.06cm]current page.north west){\color{Muted}\fontsize{7.7}{9.3}\selectfont 4/17};
\end{tikzpicture}\end{frame}
\begin{lrbox}{\TitleBox}\begin{minipage}[t]{18.2cm}\raggedright\sloppy\color{Ink}\fontsize{20.7}{24}\selectfont \bfseries\color{Navy}La matrice delle distanze\end{minipage}\end{lrbox}
\begin{lrbox}{\BodyBox}\begin{minipage}[t]{10.25cm}\raggedright\sloppy\color{Ink}\fontsize{12.4}{16.4}\selectfont La matrice D riporta la distanza fra ogni coppia di punti e conserva una diagonale nulla, poiché ogni posizione coincide con sé stessa. La simmetria deriva dalla metrica euclidea, mentre l’assenza di duplicati rende distinti i siti. Questi controlli intercettano errori di costruzione prima che le distanze entrino nei pesi del test e dell’interpolazione. Il significato geologico delle posizioni resta comunque legato alla provenienza dichiarata, che in questo esercizio corrisponde a un dominio sintetico.\end{minipage}\end{lrbox}
\begin{lrbox}{\SideBox}\begin{minipage}[t]{7.45cm}\raggedright\sloppy\color{Ink}\fontsize{12.2}{16.1}\selectfont \par\vspace{.25cm}\lstinputlisting[language=R]{estratti/P06/coordinate.R}\end{minipage}\end{lrbox}
\begin{lrbox}{\RecallBox}\begin{minipage}[t]{16.45cm}\raggedright\sloppy\color{Ink}\fontsize{10.1}{13.2}\selectfont Le proprietà della matrice verificano la geometria del calcolo e precedono l’interpretazione del campo.\end{minipage}\end{lrbox}
\MeasureBox{5-TitleBox}{\TitleBox}
\MeasureBox{5-BodyBox}{\BodyBox}
\MeasureBox{5-SideBox}{\SideBox}
\MeasureBox{5-RecallBox}{\RecallBox}
\begin{frame}[plain]\hypertarget{slide5}{}
\bookmark[level=1,page=5]{La matrice delle distanze}
\begin{tikzpicture}[remember picture,overlay]\node[anchor=north east,inner sep=0] at ([xshift=19.05cm,yshift=-0.16cm]current page.north west){\color{Classical}\fontsize{8}{10}\selectfont\hyperlink{slide2}{Indice}};
\node[anchor=north west,inner sep=0] at ([xshift=0.95cm,yshift=-0.62cm]current page.north west){\usebox{\TitleBox}};
\draw[Classical,line width=1.3pt] ([xshift=.95cm,yshift=-1.68cm]current page.north west)--([xshift=2.25cm,yshift=-1.68cm]current page.north west);\node[anchor=north west,inner sep=0] at ([xshift=0.95cm,yshift=-2.12cm]current page.north west){\usebox{\BodyBox}};
\node[anchor=north west,inner sep=0] at ([xshift=11.6cm,yshift=-2.1cm]current page.north west){\color{Quantum}\bfseries\fontsize{11.5}{14}\selectfont Implementazione in R};
\node[anchor=north west,inner sep=0] at ([xshift=11.6cm,yshift=-2.75cm]current page.north west){\usebox{\SideBox}};
\draw[Rule,line width=.5pt] ([xshift=.95cm,yshift=-10.43cm]current page.north west)--([xshift=19.05cm,yshift=-10.43cm]current page.north west);\node[anchor=north west,inner sep=0] at ([xshift=0.95cm,yshift=-10.66cm]current page.north west){\color{Classical}\bfseries\fontsize{9.4}{12}\selectfont Richiami};
\node[anchor=north west,inner sep=0] at ([xshift=2.6cm,yshift=-10.64cm]current page.north west){\usebox{\RecallBox}};
\draw[Rule,line width=.35pt] ([xshift=.95cm,yshift=-11.92cm]current page.north west)--([xshift=19.05cm,yshift=-11.92cm]current page.north west);\node[anchor=north west,inner sep=0] at ([xshift=0.95cm,yshift=-12.06cm]current page.north west){\color{Muted}\fontsize{7.7}{9.3}\selectfont P06 · Dai campioni alla validazione spaziale};
\node[anchor=north east,inner sep=0] at ([xshift=17.65cm,yshift=-12.06cm]current page.north west){\color{Classical}\fontsize{7.3}{9.1}\selectfont };
\node[anchor=north east,inner sep=0] at ([xshift=19.05cm,yshift=-12.06cm]current page.north west){\color{Muted}\fontsize{7.7}{9.3}\selectfont 5/17};
\end{tikzpicture}\end{frame}
\begin{lrbox}{\TitleBox}\begin{minipage}[t]{18.2cm}\raggedright\sloppy\color{Ink}\fontsize{20.7}{24}\selectfont \bfseries\color{Navy}La carta conserva il supporto puntuale\end{minipage}\end{lrbox}
\begin{lrbox}{\BodyBox}\begin{minipage}[t]{10.25cm}\raggedright\sloppy\color{Ink}\fontsize{12.4}{16.4}\selectfont Ogni simbolo rappresenta il valore attribuito a un sito del reticolo, mentre il colore segue una scala comune di concentrazione in unità arbitrarie. Il rapporto unitario fra gli assi conserva le distanze e rende riconoscibile la disposizione dei punti. Gli spazi fra i simboli rimangono privi di un valore osservato: un’eventuale superficie continua richiederebbe un ulteriore modello di interpolazione. La carta descrive quindi la struttura disponibile senza attribuire automaticamente la misura puntuale a tutto il suo intorno.\end{minipage}\end{lrbox}
\begin{lrbox}{\SideBox}\begin{minipage}[t]{7.45cm}\raggedright\sloppy\color{Ink}\fontsize{12.2}{16.1}\selectfont \includegraphics[width=7.45cm]{figure/P06/F01.pdf}\par\vspace{.16cm}{\fontsize{9.5}{12.6}\selectfont Sessantaquattro punti sintetici su un piano locale metrico. Colori: concentrazione in unità arbitrarie.}\end{minipage}\end{lrbox}
\begin{lrbox}{\RecallBox}\begin{minipage}[t]{16.45cm}\raggedright\sloppy\color{Ink}\fontsize{10.1}{13.2}\selectfont Supporto della misura e geometria della rappresentazione devono rimanere riconoscibili nella carta.\end{minipage}\end{lrbox}
\MeasureBox{6-TitleBox}{\TitleBox}
\MeasureBox{6-BodyBox}{\BodyBox}
\MeasureBox{6-SideBox}{\SideBox}
\MeasureBox{6-RecallBox}{\RecallBox}
\begin{frame}[plain]\hypertarget{slide6}{}
\bookmark[level=1,page=6]{La carta conserva il supporto puntuale}
\begin{tikzpicture}[remember picture,overlay]\node[anchor=north east,inner sep=0] at ([xshift=19.05cm,yshift=-0.16cm]current page.north west){\color{Classical}\fontsize{8}{10}\selectfont\hyperlink{slide2}{Indice}};
\node[anchor=north west,inner sep=0] at ([xshift=0.95cm,yshift=-0.62cm]current page.north west){\usebox{\TitleBox}};
\draw[Classical,line width=1.3pt] ([xshift=.95cm,yshift=-1.68cm]current page.north west)--([xshift=2.25cm,yshift=-1.68cm]current page.north west);\node[anchor=north west,inner sep=0] at ([xshift=0.95cm,yshift=-2.12cm]current page.north west){\usebox{\BodyBox}};
\node[anchor=north west,inner sep=0] at ([xshift=11.6cm,yshift=-2.1cm]current page.north west){\color{Quantum}\bfseries\fontsize{11.5}{14}\selectfont Evidenza grafica};
\node[anchor=north west,inner sep=0] at ([xshift=11.6cm,yshift=-2.75cm]current page.north west){\usebox{\SideBox}};
\draw[Rule,line width=.5pt] ([xshift=.95cm,yshift=-10.43cm]current page.north west)--([xshift=19.05cm,yshift=-10.43cm]current page.north west);\node[anchor=north west,inner sep=0] at ([xshift=0.95cm,yshift=-10.66cm]current page.north west){\color{Classical}\bfseries\fontsize{9.4}{12}\selectfont Richiami};
\node[anchor=north west,inner sep=0] at ([xshift=2.6cm,yshift=-10.64cm]current page.north west){\usebox{\RecallBox}};
\draw[Rule,line width=.35pt] ([xshift=.95cm,yshift=-11.92cm]current page.north west)--([xshift=19.05cm,yshift=-11.92cm]current page.north west);\node[anchor=north west,inner sep=0] at ([xshift=0.95cm,yshift=-12.06cm]current page.north west){\color{Muted}\fontsize{7.7}{9.3}\selectfont P06 · Dai campioni alla validazione spaziale};
\node[anchor=north east,inner sep=0] at ([xshift=17.65cm,yshift=-12.06cm]current page.north west){\color{Classical}\fontsize{7.3}{9.1}\selectfont };
\node[anchor=north east,inner sep=0] at ([xshift=19.05cm,yshift=-12.06cm]current page.north west){\color{Muted}\fontsize{7.7}{9.3}\selectfont 6/17};
\end{tikzpicture}\end{frame}
\begin{lrbox}{\TitleBox}\begin{minipage}[t]{18.2cm}\raggedright\sloppy\color{Ink}\fontsize{20.7}{24}\selectfont \bfseries\color{Navy}La soglia di distanza definisce una rete di relazioni\end{minipage}\end{lrbox}
\begin{lrbox}{\BodyBox}\begin{minipage}[t]{10.25cm}\raggedright\sloppy\color{Ink}\fontsize{12.4}{16.4}\selectfont Due siti sono considerati vicini quando la loro distanza è positiva e non supera 260 metri. Su questo reticolo la regola include i contatti orizzontali e verticali, escludendo quelli diagonali. Si ottengono 112 legami non orientati, che compaiono due volte nella matrice simmetrica dei pesi. La soglia definisce la scala di associazione esaminata e non rappresenta una proprietà universale del processo. Cambiarla modificherebbe il riferimento geometrico dell’indice e richiederebbe una sensibilità dichiarata.\par\vspace{.28cm}{\centering\fontsize{13}{17}\selectfont\(\displaystyle\begin{gathered}w_{ij}=\mathbf1\{0<d_{ij}\leq260\},\qquad S_0=\sum_i\sum_jw_{ij}=224\end{gathered}\)\par}\end{minipage}\end{lrbox}
\begin{lrbox}{\SideBox}\begin{minipage}[t]{7.45cm}\raggedright\sloppy\color{Ink}\fontsize{12.2}{16.1}\selectfont \includegraphics[width=7.45cm]{figure/P06/F02.pdf}\par\vspace{.16cm}{\fontsize{9.5}{12.6}\selectfont $d_{ij}$: distanza in metri. 112 legami non orientati; $S_0=224$ contributi.}\end{minipage}\end{lrbox}
\begin{lrbox}{\RecallBox}\begin{minipage}[t]{16.45cm}\raggedright\sloppy\color{Ink}\fontsize{10.1}{13.2}\selectfont L’associazione spaziale viene definita rispetto a una specifica struttura di vicinato.\end{minipage}\end{lrbox}
\MeasureBox{7-TitleBox}{\TitleBox}
\MeasureBox{7-BodyBox}{\BodyBox}
\MeasureBox{7-SideBox}{\SideBox}
\MeasureBox{7-RecallBox}{\RecallBox}
\begin{frame}[plain]\hypertarget{slide7}{}
\bookmark[level=0,page=7]{Vicinato e riferimento di permutazione}
\bookmark[level=1,page=7]{La soglia di distanza definisce una rete di relazioni}
\begin{tikzpicture}[remember picture,overlay]\node[anchor=north east,inner sep=0] at ([xshift=19.05cm,yshift=-0.16cm]current page.north west){\color{Classical}\fontsize{8}{10}\selectfont\hyperlink{slide2}{Indice}};
\node[anchor=north west,inner sep=0] at ([xshift=0.95cm,yshift=-0.62cm]current page.north west){\usebox{\TitleBox}};
\draw[Classical,line width=1.3pt] ([xshift=.95cm,yshift=-1.68cm]current page.north west)--([xshift=2.25cm,yshift=-1.68cm]current page.north west);\node[anchor=north west,inner sep=0] at ([xshift=0.95cm,yshift=-2.12cm]current page.north west){\usebox{\BodyBox}};
\node[anchor=north west,inner sep=0] at ([xshift=11.6cm,yshift=-2.1cm]current page.north west){\color{Quantum}\bfseries\fontsize{11.5}{14}\selectfont Evidenza grafica};
\node[anchor=north west,inner sep=0] at ([xshift=11.6cm,yshift=-2.75cm]current page.north west){\usebox{\SideBox}};
\draw[Rule,line width=.5pt] ([xshift=.95cm,yshift=-10.43cm]current page.north west)--([xshift=19.05cm,yshift=-10.43cm]current page.north west);\node[anchor=north west,inner sep=0] at ([xshift=0.95cm,yshift=-10.66cm]current page.north west){\color{Classical}\bfseries\fontsize{9.4}{12}\selectfont Richiami};
\node[anchor=north west,inner sep=0] at ([xshift=2.6cm,yshift=-10.64cm]current page.north west){\usebox{\RecallBox}};
\draw[Rule,line width=.35pt] ([xshift=.95cm,yshift=-11.92cm]current page.north west)--([xshift=19.05cm,yshift=-11.92cm]current page.north west);\node[anchor=north west,inner sep=0] at ([xshift=0.95cm,yshift=-12.06cm]current page.north west){\color{Muted}\fontsize{7.7}{9.3}\selectfont P06 · Dai campioni alla validazione spaziale};
\node[anchor=north east,inner sep=0] at ([xshift=17.65cm,yshift=-12.06cm]current page.north west){\color{Classical}\fontsize{7.3}{9.1}\selectfont };
\node[anchor=north east,inner sep=0] at ([xshift=19.05cm,yshift=-12.06cm]current page.north west){\color{Muted}\fontsize{7.7}{9.3}\selectfont 7/17};
\end{tikzpicture}\end{frame}
\begin{lrbox}{\TitleBox}\begin{minipage}[t]{18.2cm}\raggedright\sloppy\color{Ink}\fontsize{20.7}{24}\selectfont \bfseries\color{Navy}L’indice confronta prodotti fra scarti vicini\end{minipage}\end{lrbox}
\begin{lrbox}{\BodyBox}\begin{minipage}[t]{10.25cm}\raggedright\sloppy\color{Ink}\fontsize{12.4}{16.4}\selectfont Centrando le concentrazioni rispetto alla media del reticolo, il prodotto fra due scarti è positivo quando entrambi si trovano dalla stessa parte della media. L’indice di Moran riunisce questi prodotti sui legami selezionati e li normalizza mediante dispersione e somma dei pesi. Nel caso osservato vale circa 0,71953. Il valore indica un’associazione positiva rispetto al vicinato binario adottato, mentre la sua interpretazione inferenziale richiede la costruzione di un riferimento coerente con le ipotesi di scambiabilità.\par\vspace{.28cm}{\centering\fontsize{13}{17}\selectfont\(\displaystyle\begin{gathered}v_i=z_i-\overline z,\qquad I=\frac{n}{S_0}\frac{\sum_i\sum_jw_{ij}v_iv_j}{\sum_i v_i^2}\end{gathered}\)\par}\end{minipage}\end{lrbox}
\begin{lrbox}{\SideBox}\begin{minipage}[t]{7.45cm}\raggedright\sloppy\color{Ink}\fontsize{12.2}{16.1}\selectfont \input{tikz/P06/D02.tex}\par\vspace{.16cm}{\fontsize{9.5}{12.6}\selectfont $z_i$: concentrazione; $\bar z$: media; n = 64; $w_{ij}=1$ per $0<d_{ij}\leq260$ m, 0 altrimenti; $S_0=224$.}\end{minipage}\end{lrbox}
\begin{lrbox}{\RecallBox}\begin{minipage}[t]{16.45cm}\raggedright\sloppy\color{Ink}\fontsize{10.1}{13.2}\selectfont Segno e grandezza dell’indice dipendono da scarti, dispersione e pesi specificati.\end{minipage}\end{lrbox}
\MeasureBox{8-TitleBox}{\TitleBox}
\MeasureBox{8-BodyBox}{\BodyBox}
\MeasureBox{8-SideBox}{\SideBox}
\MeasureBox{8-RecallBox}{\RecallBox}
\begin{frame}[plain]\hypertarget{slide8}{}
\bookmark[level=1,page=8]{L’indice confronta prodotti fra scarti vicini}
\begin{tikzpicture}[remember picture,overlay]\node[anchor=north east,inner sep=0] at ([xshift=19.05cm,yshift=-0.16cm]current page.north west){\color{Classical}\fontsize{8}{10}\selectfont\hyperlink{slide2}{Indice}};
\node[anchor=north west,inner sep=0] at ([xshift=0.95cm,yshift=-0.62cm]current page.north west){\usebox{\TitleBox}};
\draw[Classical,line width=1.3pt] ([xshift=.95cm,yshift=-1.68cm]current page.north west)--([xshift=2.25cm,yshift=-1.68cm]current page.north west);\node[anchor=north west,inner sep=0] at ([xshift=0.95cm,yshift=-2.12cm]current page.north west){\usebox{\BodyBox}};
\node[anchor=north west,inner sep=0] at ([xshift=11.6cm,yshift=-2.1cm]current page.north west){\color{Quantum}\bfseries\fontsize{11.5}{14}\selectfont Lettura del diagramma};
\node[anchor=north west,inner sep=0] at ([xshift=11.6cm,yshift=-2.75cm]current page.north west){\usebox{\SideBox}};
\draw[Rule,line width=.5pt] ([xshift=.95cm,yshift=-10.43cm]current page.north west)--([xshift=19.05cm,yshift=-10.43cm]current page.north west);\node[anchor=north west,inner sep=0] at ([xshift=0.95cm,yshift=-10.66cm]current page.north west){\color{Classical}\bfseries\fontsize{9.4}{12}\selectfont Richiami};
\node[anchor=north west,inner sep=0] at ([xshift=2.6cm,yshift=-10.64cm]current page.north west){\usebox{\RecallBox}};
\draw[Rule,line width=.35pt] ([xshift=.95cm,yshift=-11.92cm]current page.north west)--([xshift=19.05cm,yshift=-11.92cm]current page.north west);\node[anchor=north west,inner sep=0] at ([xshift=0.95cm,yshift=-12.06cm]current page.north west){\color{Muted}\fontsize{7.7}{9.3}\selectfont P06 · Dai campioni alla validazione spaziale};
\node[anchor=north east,inner sep=0] at ([xshift=17.65cm,yshift=-12.06cm]current page.north west){\color{Classical}\fontsize{7.3}{9.1}\selectfont };
\node[anchor=north east,inner sep=0] at ([xshift=19.05cm,yshift=-12.06cm]current page.north west){\color{Muted}\fontsize{7.7}{9.3}\selectfont 8/17};
\end{tikzpicture}\end{frame}
\begin{lrbox}{\TitleBox}\begin{minipage}[t]{18.2cm}\raggedright\sloppy\color{Ink}\fontsize{20.7}{24}\selectfont \bfseries\color{Navy}Le permutazioni conservano siti e valori marginali\end{minipage}\end{lrbox}
\begin{lrbox}{\BodyBox}\begin{minipage}[t]{10.25cm}\raggedright\sloppy\color{Ink}\fontsize{12.4}{16.4}\selectfont Il riferimento riassegna le concentrazioni alle posizioni fisse e ricalcola l’indice, conservando insieme dei valori e geometria. Questa operazione è pertinente sotto un’ipotesi di scambiabilità delle assegnazioni, che un gradiente spaziale sistematico può rendere scientificamente discutibile come modello nullo. Nel progetto vengono conservati 999 piani di permutazione, così che la distribuzione possa essere ricostruita anche senza ripetere il generatore casuale. L’esito riguarda il riferimento scelto e non identifica il meccanismo fisico che genera la struttura.\end{minipage}\end{lrbox}
\begin{lrbox}{\SideBox}\begin{minipage}[t]{7.45cm}\raggedright\sloppy\color{Ink}\fontsize{12.2}{16.1}\selectfont \input{tikz/P06/D03.tex}\end{minipage}\end{lrbox}
\begin{lrbox}{\RecallBox}\begin{minipage}[t]{16.45cm}\raggedright\sloppy\color{Ink}\fontsize{10.1}{13.2}\selectfont La distribuzione di permutazione deriva da un’ipotesi sulle assegnazioni ammissibili.\end{minipage}\end{lrbox}
\MeasureBox{9-TitleBox}{\TitleBox}
\MeasureBox{9-BodyBox}{\BodyBox}
\MeasureBox{9-SideBox}{\SideBox}
\MeasureBox{9-RecallBox}{\RecallBox}
\begin{frame}[plain]\hypertarget{slide9}{}
\bookmark[level=1,page=9]{Le permutazioni conservano siti e valori marginali}
\begin{tikzpicture}[remember picture,overlay]\node[anchor=north east,inner sep=0] at ([xshift=19.05cm,yshift=-0.16cm]current page.north west){\color{Classical}\fontsize{8}{10}\selectfont\hyperlink{slide2}{Indice}};
\node[anchor=north west,inner sep=0] at ([xshift=0.95cm,yshift=-0.62cm]current page.north west){\usebox{\TitleBox}};
\draw[Classical,line width=1.3pt] ([xshift=.95cm,yshift=-1.68cm]current page.north west)--([xshift=2.25cm,yshift=-1.68cm]current page.north west);\node[anchor=north west,inner sep=0] at ([xshift=0.95cm,yshift=-2.12cm]current page.north west){\usebox{\BodyBox}};
\node[anchor=north west,inner sep=0] at ([xshift=11.6cm,yshift=-2.1cm]current page.north west){\color{Quantum}\bfseries\fontsize{11.5}{14}\selectfont Lettura del diagramma};
\node[anchor=north west,inner sep=0] at ([xshift=11.6cm,yshift=-2.75cm]current page.north west){\usebox{\SideBox}};
\draw[Rule,line width=.5pt] ([xshift=.95cm,yshift=-10.43cm]current page.north west)--([xshift=19.05cm,yshift=-10.43cm]current page.north west);\node[anchor=north west,inner sep=0] at ([xshift=0.95cm,yshift=-10.66cm]current page.north west){\color{Classical}\bfseries\fontsize{9.4}{12}\selectfont Richiami};
\node[anchor=north west,inner sep=0] at ([xshift=2.6cm,yshift=-10.64cm]current page.north west){\usebox{\RecallBox}};
\draw[Rule,line width=.35pt] ([xshift=.95cm,yshift=-11.92cm]current page.north west)--([xshift=19.05cm,yshift=-11.92cm]current page.north west);\node[anchor=north west,inner sep=0] at ([xshift=0.95cm,yshift=-12.06cm]current page.north west){\color{Muted}\fontsize{7.7}{9.3}\selectfont P06 · Dai campioni alla validazione spaziale};
\node[anchor=north east,inner sep=0] at ([xshift=17.65cm,yshift=-12.06cm]current page.north west){\color{Classical}\fontsize{7.3}{9.1}\selectfont };
\node[anchor=north east,inner sep=0] at ([xshift=19.05cm,yshift=-12.06cm]current page.north west){\color{Muted}\fontsize{7.7}{9.3}\selectfont 9/17};
\end{tikzpicture}\end{frame}
\begin{lrbox}{\TitleBox}\begin{minipage}[t]{18.2cm}\raggedright\sloppy\color{Ink}\fontsize{20.7}{24}\selectfont \bfseries\color{Navy}Numero di permutazioni e risoluzione\end{minipage}\end{lrbox}
\begin{lrbox}{\BodyBox}\begin{minipage}[t]{10.25cm}\raggedright\sloppy\color{Ink}\fontsize{12.4}{16.4}\selectfont Nessuno dei 999 indici permutati raggiunge quello osservato nella coda positiva. Con la correzione più uno, il p-value Monte Carlo è 0,001, il minimo rappresentabile da questa costruzione. L’assenza di superamenti nel campione simulato non implica una probabilità nulla nel riferimento completo. Aumentare le repliche affina la risoluzione numerica, lasciando immutata l’informazione osservativa del reticolo. La conclusione riguarda l’incompatibilità con l’assegnazione scambiabile considerata e dipende dalla plausibilità scientifica di tale ipotesi.\par\vspace{.28cm}{\centering\fontsize{13}{17}\selectfont\(\displaystyle\begin{gathered}\widehat p=\frac{1+\sum_{b=1}^{999}\mathbf1\{I_b\geq I_{obs}\}}{1000}=0{,}001\end{gathered}\)\par}\end{minipage}\end{lrbox}
\begin{lrbox}{\SideBox}\begin{minipage}[t]{7.45cm}\raggedright\sloppy\color{Ink}\fontsize{12.2}{16.1}\selectfont \includegraphics[width=7.45cm]{figure/P06/F03.pdf}\par\vspace{.16cm}{\fontsize{9.5}{12.6}\selectfont $I_b$: indice nella permutazione b; $I_{obs}=0,71953$: indice osservato, in viola. Istogramma: 999 valori permutati.}\end{minipage}\end{lrbox}
\begin{lrbox}{\RecallBox}\begin{minipage}[t]{16.45cm}\raggedright\sloppy\color{Ink}\fontsize{10.1}{13.2}\selectfont La precisione Monte Carlo e la plausibilità scientifica del riferimento sono componenti distinte del test.\end{minipage}\end{lrbox}
\MeasureBox{10-TitleBox}{\TitleBox}
\MeasureBox{10-BodyBox}{\BodyBox}
\MeasureBox{10-SideBox}{\SideBox}
\MeasureBox{10-RecallBox}{\RecallBox}
\begin{frame}[plain]\hypertarget{slide10}{}
\bookmark[level=1,page=10]{Numero di permutazioni e risoluzione}
\begin{tikzpicture}[remember picture,overlay]\node[anchor=north east,inner sep=0] at ([xshift=19.05cm,yshift=-0.16cm]current page.north west){\color{Classical}\fontsize{8}{10}\selectfont\hyperlink{slide2}{Indice}};
\node[anchor=north west,inner sep=0] at ([xshift=0.95cm,yshift=-0.62cm]current page.north west){\usebox{\TitleBox}};
\draw[Classical,line width=1.3pt] ([xshift=.95cm,yshift=-1.68cm]current page.north west)--([xshift=2.25cm,yshift=-1.68cm]current page.north west);\node[anchor=north west,inner sep=0] at ([xshift=0.95cm,yshift=-2.12cm]current page.north west){\usebox{\BodyBox}};
\node[anchor=north west,inner sep=0] at ([xshift=11.6cm,yshift=-2.1cm]current page.north west){\color{Quantum}\bfseries\fontsize{11.5}{14}\selectfont Evidenza grafica};
\node[anchor=north west,inner sep=0] at ([xshift=11.6cm,yshift=-2.75cm]current page.north west){\usebox{\SideBox}};
\draw[Rule,line width=.5pt] ([xshift=.95cm,yshift=-10.43cm]current page.north west)--([xshift=19.05cm,yshift=-10.43cm]current page.north west);\node[anchor=north west,inner sep=0] at ([xshift=0.95cm,yshift=-10.66cm]current page.north west){\color{Classical}\bfseries\fontsize{9.4}{12}\selectfont Richiami};
\node[anchor=north west,inner sep=0] at ([xshift=2.6cm,yshift=-10.64cm]current page.north west){\usebox{\RecallBox}};
\draw[Rule,line width=.35pt] ([xshift=.95cm,yshift=-11.92cm]current page.north west)--([xshift=19.05cm,yshift=-11.92cm]current page.north west);\node[anchor=north west,inner sep=0] at ([xshift=0.95cm,yshift=-12.06cm]current page.north west){\color{Muted}\fontsize{7.7}{9.3}\selectfont P06 · Dai campioni alla validazione spaziale};
\node[anchor=north east,inner sep=0] at ([xshift=17.65cm,yshift=-12.06cm]current page.north west){\color{Classical}\fontsize{7.3}{9.1}\selectfont };
\node[anchor=north east,inner sep=0] at ([xshift=19.05cm,yshift=-12.06cm]current page.north west){\color{Muted}\fontsize{7.7}{9.3}\selectfont 10/17};
\end{tikzpicture}\end{frame}
\begin{lrbox}{\TitleBox}\begin{minipage}[t]{18.2cm}\raggedright\sloppy\color{Ink}\fontsize{20.7}{24}\selectfont \bfseries\color{Navy}La geometria del settore escluso\end{minipage}\end{lrbox}
\begin{lrbox}{\BodyBox}\begin{minipage}[t]{10.25cm}\raggedright\sloppy\color{Ink}\fontsize{12.4}{16.4}\selectfont L’adattamento utilizza i quaranta siti con coordinata x inferiore a 1250 metri, mentre i ventiquattro siti rimanenti costituiscono il settore orientale di verifica. Il valore massimo di x nel train è mille metri, per cui il test oltrepassa la copertura osservata lungo quella direzione. Questa separazione rappresenta un trasferimento spaziale e richiede estrapolazione geografica rispetto ai siti disponibili. Una divisione casuale avrebbe spesso vicini di adattamento attorno ai punti di verifica e risponderebbe a una domanda differente.\end{minipage}\end{lrbox}
\begin{lrbox}{\SideBox}\begin{minipage}[t]{7.45cm}\raggedright\sloppy\color{Ink}\fontsize{12.2}{16.1}\selectfont \includegraphics[width=7.45cm]{figure/P06/F04.pdf}\par\vspace{.16cm}{\fontsize{9.5}{12.6}\selectfont Verde: quaranta siti di adattamento; viola: ventiquattro siti orientali separati per la verifica.}\end{minipage}\end{lrbox}
\begin{lrbox}{\RecallBox}\begin{minipage}[t]{16.45cm}\raggedright\sloppy\color{Ink}\fontsize{10.1}{13.2}\selectfont La geometria della separazione determina il tipo di generalizzazione che viene valutato.\end{minipage}\end{lrbox}
\MeasureBox{11-TitleBox}{\TitleBox}
\MeasureBox{11-BodyBox}{\BodyBox}
\MeasureBox{11-SideBox}{\SideBox}
\MeasureBox{11-RecallBox}{\RecallBox}
\begin{frame}[plain]\hypertarget{slide11}{}
\bookmark[level=0,page=11]{Previsione e verifica nel settore orientale}
\bookmark[level=1,page=11]{La geometria del settore escluso}
\begin{tikzpicture}[remember picture,overlay]\node[anchor=north east,inner sep=0] at ([xshift=19.05cm,yshift=-0.16cm]current page.north west){\color{Classical}\fontsize{8}{10}\selectfont\hyperlink{slide2}{Indice}};
\node[anchor=north west,inner sep=0] at ([xshift=0.95cm,yshift=-0.62cm]current page.north west){\usebox{\TitleBox}};
\draw[Classical,line width=1.3pt] ([xshift=.95cm,yshift=-1.68cm]current page.north west)--([xshift=2.25cm,yshift=-1.68cm]current page.north west);\node[anchor=north west,inner sep=0] at ([xshift=0.95cm,yshift=-2.12cm]current page.north west){\usebox{\BodyBox}};
\node[anchor=north west,inner sep=0] at ([xshift=11.6cm,yshift=-2.1cm]current page.north west){\color{Quantum}\bfseries\fontsize{11.5}{14}\selectfont Evidenza grafica};
\node[anchor=north west,inner sep=0] at ([xshift=11.6cm,yshift=-2.75cm]current page.north west){\usebox{\SideBox}};
\draw[Rule,line width=.5pt] ([xshift=.95cm,yshift=-10.43cm]current page.north west)--([xshift=19.05cm,yshift=-10.43cm]current page.north west);\node[anchor=north west,inner sep=0] at ([xshift=0.95cm,yshift=-10.66cm]current page.north west){\color{Classical}\bfseries\fontsize{9.4}{12}\selectfont Richiami};
\node[anchor=north west,inner sep=0] at ([xshift=2.6cm,yshift=-10.64cm]current page.north west){\usebox{\RecallBox}};
\draw[Rule,line width=.35pt] ([xshift=.95cm,yshift=-11.92cm]current page.north west)--([xshift=19.05cm,yshift=-11.92cm]current page.north west);\node[anchor=north west,inner sep=0] at ([xshift=0.95cm,yshift=-12.06cm]current page.north west){\color{Muted}\fontsize{7.7}{9.3}\selectfont P06 · Dai campioni alla validazione spaziale};
\node[anchor=north east,inner sep=0] at ([xshift=17.65cm,yshift=-12.06cm]current page.north west){\color{Classical}\fontsize{7.3}{9.1}\selectfont };
\node[anchor=north east,inner sep=0] at ([xshift=19.05cm,yshift=-12.06cm]current page.north west){\color{Muted}\fontsize{7.7}{9.3}\selectfont 11/17};
\end{tikzpicture}\end{frame}
\begin{lrbox}{\TitleBox}\begin{minipage}[t]{18.2cm}\raggedright\sloppy\color{Ink}\fontsize{20.7}{24}\selectfont \bfseries\color{Navy}L’IDW costruisce una combinazione convessa locale\end{minipage}\end{lrbox}
\begin{lrbox}{\BodyBox}\begin{minipage}[t]{10.25cm}\raggedright\sloppy\color{Ink}\fontsize{12.4}{16.4}\selectfont La previsione usa pesi inversamente proporzionali al quadrato della distanza dai soli siti di adattamento, normalizzati affinché la loro somma sia uno. La previsione appartiene quindi all’intervallo dei valori disponibili nel train. Quando un gradiente prosegue oltre il settore osservato, questo vincolo può impedire di riprodurre valori più alti presenti nella fascia esclusa. Il calcolo del peso dipende soltanto dalla posizione del punto da prevedere e dai dati ammessi all’adattamento.\par\vspace{.28cm}{\centering\fontsize{13}{17}\selectfont\(\displaystyle\begin{gathered}\lambda_i(s)=\frac{d(s,s_i)^{-2}}{\sum_{j\in train}d(s,s_j)^{-2}},\qquad \widehat z(s)=\sum_{i\in train}\lambda_i(s)z_i\end{gathered}\)\par}\end{minipage}\end{lrbox}
\begin{lrbox}{\SideBox}\begin{minipage}[t]{7.45cm}\raggedright\sloppy\color{Ink}\fontsize{12.2}{16.1}\selectfont \input{tikz/P06/D04.tex}\par\vspace{.16cm}{\fontsize{9.5}{12.6}\selectfont s: posizione da prevedere; $s_i,z_i$: sito e valore di adattamento; d: distanza in metri, qui positiva. $\lambda_i$: peso normalizzato.}\end{minipage}\end{lrbox}
\begin{lrbox}{\RecallBox}\begin{minipage}[t]{16.45cm}\raggedright\sloppy\color{Ink}\fontsize{10.1}{13.2}\selectfont La positività dei pesi limita la previsione ai valori osservati nei siti utilizzati.\end{minipage}\end{lrbox}
\MeasureBox{12-TitleBox}{\TitleBox}
\MeasureBox{12-BodyBox}{\BodyBox}
\MeasureBox{12-SideBox}{\SideBox}
\MeasureBox{12-RecallBox}{\RecallBox}
\begin{frame}[plain]\hypertarget{slide12}{}
\bookmark[level=1,page=12]{L’IDW costruisce una combinazione convessa locale}
\begin{tikzpicture}[remember picture,overlay]\node[anchor=north east,inner sep=0] at ([xshift=19.05cm,yshift=-0.16cm]current page.north west){\color{Classical}\fontsize{8}{10}\selectfont\hyperlink{slide2}{Indice}};
\node[anchor=north west,inner sep=0] at ([xshift=0.95cm,yshift=-0.62cm]current page.north west){\usebox{\TitleBox}};
\draw[Classical,line width=1.3pt] ([xshift=.95cm,yshift=-1.68cm]current page.north west)--([xshift=2.25cm,yshift=-1.68cm]current page.north west);\node[anchor=north west,inner sep=0] at ([xshift=0.95cm,yshift=-2.12cm]current page.north west){\usebox{\BodyBox}};
\node[anchor=north west,inner sep=0] at ([xshift=11.6cm,yshift=-2.1cm]current page.north west){\color{Quantum}\bfseries\fontsize{11.5}{14}\selectfont Lettura del diagramma};
\node[anchor=north west,inner sep=0] at ([xshift=11.6cm,yshift=-2.75cm]current page.north west){\usebox{\SideBox}};
\draw[Rule,line width=.5pt] ([xshift=.95cm,yshift=-10.43cm]current page.north west)--([xshift=19.05cm,yshift=-10.43cm]current page.north west);\node[anchor=north west,inner sep=0] at ([xshift=0.95cm,yshift=-10.66cm]current page.north west){\color{Classical}\bfseries\fontsize{9.4}{12}\selectfont Richiami};
\node[anchor=north west,inner sep=0] at ([xshift=2.6cm,yshift=-10.64cm]current page.north west){\usebox{\RecallBox}};
\draw[Rule,line width=.35pt] ([xshift=.95cm,yshift=-11.92cm]current page.north west)--([xshift=19.05cm,yshift=-11.92cm]current page.north west);\node[anchor=north west,inner sep=0] at ([xshift=0.95cm,yshift=-12.06cm]current page.north west){\color{Muted}\fontsize{7.7}{9.3}\selectfont P06 · Dai campioni alla validazione spaziale};
\node[anchor=north east,inner sep=0] at ([xshift=17.65cm,yshift=-12.06cm]current page.north west){\color{Classical}\fontsize{7.3}{9.1}\selectfont };
\node[anchor=north east,inner sep=0] at ([xshift=19.05cm,yshift=-12.06cm]current page.north west){\color{Muted}\fontsize{7.7}{9.3}\selectfont 12/17};
\end{tikzpicture}\end{frame}
\begin{lrbox}{\TitleBox}\begin{minipage}[t]{18.2cm}\raggedright\sloppy\color{Ink}\fontsize{20.7}{24}\selectfont \bfseries\color{Navy}Il codice esclude le risposte del settore di verifica\end{minipage}\end{lrbox}
\begin{lrbox}{\BodyBox}\begin{minipage}[t]{10.25cm}\raggedright\sloppy\color{Ink}\fontsize{12.4}{16.4}\selectfont Il frammento percorre i siti orientali e costruisce per ciascuno le distanze dai soli quaranta siti di adattamento. Le concentrazioni orientali entrano soltanto nel calcolo finale degli errori, dopo che le previsioni sono state prodotte. La separazione viene verificata insieme alla finitezza dei risultati e all’appartenenza delle previsioni all’intervallo del train. Una funzione generale gestisce inoltre il caso di distanza nulla, restituendo il valore del sito coincidente; nella fascia esclusa questo caso non si presenta.\end{minipage}\end{lrbox}
\begin{lrbox}{\SideBox}\begin{minipage}[t]{7.45cm}\raggedright\sloppy\color{Ink}\fontsize{12.2}{16.1}\selectfont \par\vspace{.25cm}\lstinputlisting[language=R]{estratti/P06/IDW.R}\end{minipage}\end{lrbox}
\begin{lrbox}{\RecallBox}\begin{minipage}[t]{16.45cm}\raggedright\sloppy\color{Ink}\fontsize{10.1}{13.2}\selectfont La verifica indipendente richiede che i valori da prevedere siano esclusi dalla costruzione dei pesi e della risposta stimata.\end{minipage}\end{lrbox}
\MeasureBox{13-TitleBox}{\TitleBox}
\MeasureBox{13-BodyBox}{\BodyBox}
\MeasureBox{13-SideBox}{\SideBox}
\MeasureBox{13-RecallBox}{\RecallBox}
\begin{frame}[plain]\hypertarget{slide13}{}
\bookmark[level=1,page=13]{Il codice esclude le risposte del settore di verifica}
\begin{tikzpicture}[remember picture,overlay]\node[anchor=north east,inner sep=0] at ([xshift=19.05cm,yshift=-0.16cm]current page.north west){\color{Classical}\fontsize{8}{10}\selectfont\hyperlink{slide2}{Indice}};
\node[anchor=north west,inner sep=0] at ([xshift=0.95cm,yshift=-0.62cm]current page.north west){\usebox{\TitleBox}};
\draw[Classical,line width=1.3pt] ([xshift=.95cm,yshift=-1.68cm]current page.north west)--([xshift=2.25cm,yshift=-1.68cm]current page.north west);\node[anchor=north west,inner sep=0] at ([xshift=0.95cm,yshift=-2.12cm]current page.north west){\usebox{\BodyBox}};
\node[anchor=north west,inner sep=0] at ([xshift=11.6cm,yshift=-2.1cm]current page.north west){\color{Quantum}\bfseries\fontsize{11.5}{14}\selectfont Implementazione in R};
\node[anchor=north west,inner sep=0] at ([xshift=11.6cm,yshift=-2.75cm]current page.north west){\usebox{\SideBox}};
\draw[Rule,line width=.5pt] ([xshift=.95cm,yshift=-10.43cm]current page.north west)--([xshift=19.05cm,yshift=-10.43cm]current page.north west);\node[anchor=north west,inner sep=0] at ([xshift=0.95cm,yshift=-10.66cm]current page.north west){\color{Classical}\bfseries\fontsize{9.4}{12}\selectfont Richiami};
\node[anchor=north west,inner sep=0] at ([xshift=2.6cm,yshift=-10.64cm]current page.north west){\usebox{\RecallBox}};
\draw[Rule,line width=.35pt] ([xshift=.95cm,yshift=-11.92cm]current page.north west)--([xshift=19.05cm,yshift=-11.92cm]current page.north west);\node[anchor=north west,inner sep=0] at ([xshift=0.95cm,yshift=-12.06cm]current page.north west){\color{Muted}\fontsize{7.7}{9.3}\selectfont P06 · Dai campioni alla validazione spaziale};
\node[anchor=north east,inner sep=0] at ([xshift=17.65cm,yshift=-12.06cm]current page.north west){\color{Classical}\fontsize{7.3}{9.1}\selectfont };
\node[anchor=north east,inner sep=0] at ([xshift=19.05cm,yshift=-12.06cm]current page.north west){\color{Muted}\fontsize{7.7}{9.3}\selectfont 13/17};
\end{tikzpicture}\end{frame}
\begin{lrbox}{\TitleBox}\begin{minipage}[t]{18.2cm}\raggedright\sloppy\color{Ink}\fontsize{20.7}{24}\selectfont \bfseries\color{Navy}La carta degli errori conserva la loro localizzazione\end{minipage}\end{lrbox}
\begin{lrbox}{\BodyBox}\begin{minipage}[t]{10.25cm}\raggedright\sloppy\color{Ink}\fontsize{12.4}{16.4}\selectfont Gli errori sono definiti come osservato meno previsto, per cui un valore positivo indica sottostima. La carta mostra dove la previsione si discosti maggiormente dal campo sintetico, mantenendo visibili i siti di adattamento in grigio. La media degli errori è circa 3,175 unità arbitrarie, segnalando una sottostima complessiva nel settore orientale. La localizzazione aiuta a collegare questa direzione dello scarto al gradiente e al vincolo delle combinazioni convesse, anziché ridurre la diagnosi a un’unica misura aggregata.\end{minipage}\end{lrbox}
\begin{lrbox}{\SideBox}\begin{minipage}[t]{7.45cm}\raggedright\sloppy\color{Ink}\fontsize{12.2}{16.1}\selectfont \includegraphics[width=7.45cm]{figure/P06/F05.pdf}\par\vspace{.16cm}{\fontsize{9.5}{12.6}\selectfont Errori osservato meno previsto nella fascia est. Scala divergente centrata sullo zero.}\end{minipage}\end{lrbox}
\begin{lrbox}{\RecallBox}\begin{minipage}[t]{16.45cm}\raggedright\sloppy\color{Ink}\fontsize{10.1}{13.2}\selectfont La posizione e il segno degli errori possono spiegare limiti che una metrica aggregata nasconde.\end{minipage}\end{lrbox}
\MeasureBox{14-TitleBox}{\TitleBox}
\MeasureBox{14-BodyBox}{\BodyBox}
\MeasureBox{14-SideBox}{\SideBox}
\MeasureBox{14-RecallBox}{\RecallBox}
\begin{frame}[plain]\hypertarget{slide14}{}
\bookmark[level=1,page=14]{La carta degli errori conserva la loro localizzazione}
\begin{tikzpicture}[remember picture,overlay]\node[anchor=north east,inner sep=0] at ([xshift=19.05cm,yshift=-0.16cm]current page.north west){\color{Classical}\fontsize{8}{10}\selectfont\hyperlink{slide2}{Indice}};
\node[anchor=north west,inner sep=0] at ([xshift=0.95cm,yshift=-0.62cm]current page.north west){\usebox{\TitleBox}};
\draw[Classical,line width=1.3pt] ([xshift=.95cm,yshift=-1.68cm]current page.north west)--([xshift=2.25cm,yshift=-1.68cm]current page.north west);\node[anchor=north west,inner sep=0] at ([xshift=0.95cm,yshift=-2.12cm]current page.north west){\usebox{\BodyBox}};
\node[anchor=north west,inner sep=0] at ([xshift=11.6cm,yshift=-2.1cm]current page.north west){\color{Quantum}\bfseries\fontsize{11.5}{14}\selectfont Evidenza grafica};
\node[anchor=north west,inner sep=0] at ([xshift=11.6cm,yshift=-2.75cm]current page.north west){\usebox{\SideBox}};
\draw[Rule,line width=.5pt] ([xshift=.95cm,yshift=-10.43cm]current page.north west)--([xshift=19.05cm,yshift=-10.43cm]current page.north west);\node[anchor=north west,inner sep=0] at ([xshift=0.95cm,yshift=-10.66cm]current page.north west){\color{Classical}\bfseries\fontsize{9.4}{12}\selectfont Richiami};
\node[anchor=north west,inner sep=0] at ([xshift=2.6cm,yshift=-10.64cm]current page.north west){\usebox{\RecallBox}};
\draw[Rule,line width=.35pt] ([xshift=.95cm,yshift=-11.92cm]current page.north west)--([xshift=19.05cm,yshift=-11.92cm]current page.north west);\node[anchor=north west,inner sep=0] at ([xshift=0.95cm,yshift=-12.06cm]current page.north west){\color{Muted}\fontsize{7.7}{9.3}\selectfont P06 · Dai campioni alla validazione spaziale};
\node[anchor=north east,inner sep=0] at ([xshift=17.65cm,yshift=-12.06cm]current page.north west){\color{Classical}\fontsize{7.3}{9.1}\selectfont };
\node[anchor=north east,inner sep=0] at ([xshift=19.05cm,yshift=-12.06cm]current page.north west){\color{Muted}\fontsize{7.7}{9.3}\selectfont 14/17};
\end{tikzpicture}\end{frame}
\begin{lrbox}{\TitleBox}\begin{minipage}[t]{18.2cm}\raggedright\sloppy\color{Ink}\fontsize{20.7}{24}\selectfont \bfseries\color{Navy}Gli errori vengono quantificati nella scala del campo\end{minipage}\end{lrbox}
\begin{lrbox}{\BodyBox}\begin{minipage}[t]{10.25cm}\raggedright\sloppy\color{Ink}\fontsize{12.4}{16.4}\selectfont La MAE sul settore orientale è circa 3,257 e la RMSE circa 3,808 unità arbitrarie. Il diagramma osservato–previsto adotta la stessa scala sui due assi e consente di riconoscere le sottostime rispetto alla bisettrice. Queste misure descrivono il trasferimento geografico fissato, mentre il piccolo valore del test di Moran riguarda l’associazione nel reticolo. I due risultati possono dunque coesistere senza contraddizione, poiché una struttura spaziale evidente non garantisce una buona estrapolazione da una parte del dominio.\end{minipage}\end{lrbox}
\begin{lrbox}{\SideBox}\begin{minipage}[t]{7.45cm}\raggedright\sloppy\color{Ink}\fontsize{12.2}{16.1}\selectfont \includegraphics[width=7.45cm]{figure/P06/F06.pdf}\par\vspace{.16cm}{\fontsize{9.5}{12.6}\selectfont Confronto sui ventiquattro siti esclusi. La bisettrice rappresenta previsione esatta.}\end{minipage}\end{lrbox}
\begin{lrbox}{\RecallBox}\begin{minipage}[t]{16.45cm}\raggedright\sloppy\color{Ink}\fontsize{10.1}{13.2}\selectfont L’associazione spaziale descritta e la capacità di trasferimento valutata rispondono a domande differenti.\end{minipage}\end{lrbox}
\MeasureBox{15-TitleBox}{\TitleBox}
\MeasureBox{15-BodyBox}{\BodyBox}
\MeasureBox{15-SideBox}{\SideBox}
\MeasureBox{15-RecallBox}{\RecallBox}
\begin{frame}[plain]\hypertarget{slide15}{}
\bookmark[level=1,page=15]{Gli errori vengono quantificati nella scala del campo}
\begin{tikzpicture}[remember picture,overlay]\node[anchor=north east,inner sep=0] at ([xshift=19.05cm,yshift=-0.16cm]current page.north west){\color{Classical}\fontsize{8}{10}\selectfont\hyperlink{slide2}{Indice}};
\node[anchor=north west,inner sep=0] at ([xshift=0.95cm,yshift=-0.62cm]current page.north west){\usebox{\TitleBox}};
\draw[Classical,line width=1.3pt] ([xshift=.95cm,yshift=-1.68cm]current page.north west)--([xshift=2.25cm,yshift=-1.68cm]current page.north west);\node[anchor=north west,inner sep=0] at ([xshift=0.95cm,yshift=-2.12cm]current page.north west){\usebox{\BodyBox}};
\node[anchor=north west,inner sep=0] at ([xshift=11.6cm,yshift=-2.1cm]current page.north west){\color{Quantum}\bfseries\fontsize{11.5}{14}\selectfont Evidenza grafica};
\node[anchor=north west,inner sep=0] at ([xshift=11.6cm,yshift=-2.75cm]current page.north west){\usebox{\SideBox}};
\draw[Rule,line width=.5pt] ([xshift=.95cm,yshift=-10.43cm]current page.north west)--([xshift=19.05cm,yshift=-10.43cm]current page.north west);\node[anchor=north west,inner sep=0] at ([xshift=0.95cm,yshift=-10.66cm]current page.north west){\color{Classical}\bfseries\fontsize{9.4}{12}\selectfont Richiami};
\node[anchor=north west,inner sep=0] at ([xshift=2.6cm,yshift=-10.64cm]current page.north west){\usebox{\RecallBox}};
\draw[Rule,line width=.35pt] ([xshift=.95cm,yshift=-11.92cm]current page.north west)--([xshift=19.05cm,yshift=-11.92cm]current page.north west);\node[anchor=north west,inner sep=0] at ([xshift=0.95cm,yshift=-12.06cm]current page.north west){\color{Muted}\fontsize{7.7}{9.3}\selectfont P06 · Dai campioni alla validazione spaziale};
\node[anchor=north east,inner sep=0] at ([xshift=17.65cm,yshift=-12.06cm]current page.north west){\color{Classical}\fontsize{7.3}{9.1}\selectfont };
\node[anchor=north east,inner sep=0] at ([xshift=19.05cm,yshift=-12.06cm]current page.north west){\color{Muted}\fontsize{7.7}{9.3}\selectfont 15/17};
\end{tikzpicture}\end{frame}
\begin{lrbox}{\TitleBox}\begin{minipage}[t]{18.2cm}\raggedright\sloppy\color{Ink}\fontsize{20.7}{24}\selectfont \bfseries\color{Navy}La sensibilità della potenza resta esplorativa\end{minipage}\end{lrbox}
\begin{lrbox}{\BodyBox}\begin{minipage}[t]{10.25cm}\raggedright\sloppy\color{Ink}\fontsize{12.4}{16.4}\selectfont Il confronto delle potenze uno, due e tre rende visibile come una maggiore enfasi sui siti vicini modifichi l’errore nella stessa fascia est. La regola principale rimane quella con potenza due, dichiarata prima della verifica. Se si scegliesse la potenza dopo aver confrontato questi errori, il settore orientale assumerebbe il ruolo di insieme di selezione e servirebbe una nuova area indipendente per la valutazione finale. La figura viene quindi utilizzata per comprendere il metodo e i limiti del disegno.\end{minipage}\end{lrbox}
\begin{lrbox}{\SideBox}\begin{minipage}[t]{7.45cm}\raggedright\sloppy\color{Ink}\fontsize{12.2}{16.1}\selectfont \includegraphics[width=7.45cm]{figure/P06/F07.pdf}\par\vspace{.16cm}{\fontsize{9.5}{12.6}\selectfont Sensibilità descrittiva alla potenza IDW sugli stessi siti orientali; nessuna nuova verifica indipendente.}\end{minipage}\end{lrbox}
\begin{lrbox}{\RecallBox}\begin{minipage}[t]{16.45cm}\raggedright\sloppy\color{Ink}\fontsize{10.1}{13.2}\selectfont Un campione utilizzato per scegliere un parametro cambia ruolo nel procedimento di validazione.\end{minipage}\end{lrbox}
\MeasureBox{16-TitleBox}{\TitleBox}
\MeasureBox{16-BodyBox}{\BodyBox}
\MeasureBox{16-SideBox}{\SideBox}
\MeasureBox{16-RecallBox}{\RecallBox}
\begin{frame}[plain]\hypertarget{slide16}{}
\bookmark[level=1,page=16]{La sensibilità della potenza resta esplorativa}
\begin{tikzpicture}[remember picture,overlay]\node[anchor=north east,inner sep=0] at ([xshift=19.05cm,yshift=-0.16cm]current page.north west){\color{Classical}\fontsize{8}{10}\selectfont\hyperlink{slide2}{Indice}};
\node[anchor=north west,inner sep=0] at ([xshift=0.95cm,yshift=-0.62cm]current page.north west){\usebox{\TitleBox}};
\draw[Classical,line width=1.3pt] ([xshift=.95cm,yshift=-1.68cm]current page.north west)--([xshift=2.25cm,yshift=-1.68cm]current page.north west);\node[anchor=north west,inner sep=0] at ([xshift=0.95cm,yshift=-2.12cm]current page.north west){\usebox{\BodyBox}};
\node[anchor=north west,inner sep=0] at ([xshift=11.6cm,yshift=-2.1cm]current page.north west){\color{Quantum}\bfseries\fontsize{11.5}{14}\selectfont Evidenza grafica};
\node[anchor=north west,inner sep=0] at ([xshift=11.6cm,yshift=-2.75cm]current page.north west){\usebox{\SideBox}};
\draw[Rule,line width=.5pt] ([xshift=.95cm,yshift=-10.43cm]current page.north west)--([xshift=19.05cm,yshift=-10.43cm]current page.north west);\node[anchor=north west,inner sep=0] at ([xshift=0.95cm,yshift=-10.66cm]current page.north west){\color{Classical}\bfseries\fontsize{9.4}{12}\selectfont Richiami};
\node[anchor=north west,inner sep=0] at ([xshift=2.6cm,yshift=-10.64cm]current page.north west){\usebox{\RecallBox}};
\draw[Rule,line width=.35pt] ([xshift=.95cm,yshift=-11.92cm]current page.north west)--([xshift=19.05cm,yshift=-11.92cm]current page.north west);\node[anchor=north west,inner sep=0] at ([xshift=0.95cm,yshift=-12.06cm]current page.north west){\color{Muted}\fontsize{7.7}{9.3}\selectfont P06 · Dai campioni alla validazione spaziale};
\node[anchor=north east,inner sep=0] at ([xshift=17.65cm,yshift=-12.06cm]current page.north west){\color{Classical}\fontsize{7.3}{9.1}\selectfont };
\node[anchor=north east,inner sep=0] at ([xshift=19.05cm,yshift=-12.06cm]current page.north west){\color{Muted}\fontsize{7.7}{9.3}\selectfont 16/17};
\end{tikzpicture}\end{frame}
\begin{lrbox}{\TitleBox}\begin{minipage}[t]{18.2cm}\raggedright\sloppy\color{Ink}\fontsize{20.7}{24}\selectfont \bfseries\color{Navy}Consegna: collegare geometria, test e trasferimento\end{minipage}\end{lrbox}
\begin{lrbox}{\BodyBox}\begin{minipage}[t]{10.25cm}\raggedright\sloppy\color{Ink}\fontsize{12.4}{16.4}\selectfont La relazione presenta audit metrico, carta dei punti, vicinato, indice con riferimento di permutazione, separazione geografica e mappa degli errori. L’esercizio richiede di spiegare perché un indice elevato possa accompagnarsi a sottostima nel settore orientale, richiamando gradiente e vincoli dei pesi. La conclusione conserva la natura sintetica del campo e indica quali nuove misure territoriali renderebbero più informativo il trasferimento. Codice, piani di permutazione e previsioni per sito permettono di ricostruire ciascun passaggio.\end{minipage}\end{lrbox}
\begin{lrbox}{\SideBox}\begin{minipage}[t]{7.45cm}\raggedright\sloppy\color{Ink}\fontsize{12.2}{16.1}\selectfont R/P06\_figure.R integra le attività M11\_N01, N07, N14 e N19. Il progetto utilizza il vicinato binario entro 260 metri e 999 permutazioni; i laboratori dichiarano separatamente le proprie numerosità. Le coordinate geografiche approssimate descrivono una costruzione didattica, della quale non si deve inferire una provenienza da rilievo.\end{minipage}\end{lrbox}
\begin{lrbox}{\RecallBox}\begin{minipage}[t]{16.45cm}\raggedright\sloppy\color{Ink}\fontsize{10.1}{13.2}\selectfont La relazione collega la conclusione al riferimento statistico e alla geometria effettivamente valutati.\end{minipage}\end{lrbox}
\MeasureBox{17-TitleBox}{\TitleBox}
\MeasureBox{17-BodyBox}{\BodyBox}
\MeasureBox{17-SideBox}{\SideBox}
\MeasureBox{17-RecallBox}{\RecallBox}
\begin{frame}[plain]\hypertarget{slide17}{}
\bookmark[level=1,page=17]{Consegna: collegare geometria, test e trasferimento}
\begin{tikzpicture}[remember picture,overlay]\node[anchor=north east,inner sep=0] at ([xshift=19.05cm,yshift=-0.16cm]current page.north west){\color{Classical}\fontsize{8}{10}\selectfont\hyperlink{slide2}{Indice}};
\node[anchor=north west,inner sep=0] at ([xshift=0.95cm,yshift=-0.62cm]current page.north west){\usebox{\TitleBox}};
\draw[Classical,line width=1.3pt] ([xshift=.95cm,yshift=-1.68cm]current page.north west)--([xshift=2.25cm,yshift=-1.68cm]current page.north west);\node[anchor=north west,inner sep=0] at ([xshift=0.95cm,yshift=-2.12cm]current page.north west){\usebox{\BodyBox}};
\node[anchor=north west,inner sep=0] at ([xshift=11.6cm,yshift=-2.1cm]current page.north west){\color{Quantum}\bfseries\fontsize{11.5}{14}\selectfont Sviluppo e interpretazione};
\node[anchor=north west,inner sep=0] at ([xshift=11.6cm,yshift=-2.75cm]current page.north west){\usebox{\SideBox}};
\draw[Rule,line width=.5pt] ([xshift=.95cm,yshift=-10.43cm]current page.north west)--([xshift=19.05cm,yshift=-10.43cm]current page.north west);\node[anchor=north west,inner sep=0] at ([xshift=0.95cm,yshift=-10.66cm]current page.north west){\color{Classical}\bfseries\fontsize{9.4}{12}\selectfont Richiami};
\node[anchor=north west,inner sep=0] at ([xshift=2.6cm,yshift=-10.64cm]current page.north west){\usebox{\RecallBox}};
\draw[Rule,line width=.35pt] ([xshift=.95cm,yshift=-11.92cm]current page.north west)--([xshift=19.05cm,yshift=-11.92cm]current page.north west);\node[anchor=north west,inner sep=0] at ([xshift=0.95cm,yshift=-12.06cm]current page.north west){\color{Muted}\fontsize{7.7}{9.3}\selectfont P06 · Dai campioni alla validazione spaziale};
\node[anchor=north east,inner sep=0] at ([xshift=17.65cm,yshift=-12.06cm]current page.north west){\color{Classical}\fontsize{7.3}{9.1}\selectfont Raccordo M11; documentazione sf sui sistemi di riferimento};
\node[anchor=north east,inner sep=0] at ([xshift=19.05cm,yshift=-12.06cm]current page.north west){\color{Muted}\fontsize{7.7}{9.3}\selectfont 17/17};
\end{tikzpicture}\end{frame}
\end{document}